\subsection{Two target-side exclusions for a point kernel}
\paragraph{A2Y (Proved under either stated target hypothesis).}
Let $X/k$ be smooth, projective and integral of dimension
$n\geq3$ over an algebraically closed field of characteristic
$p>0$, and let $Y\subset X$ be a smooth very ample divisor.
Restriction
\[
 \operatorname{Pic}(X)\longrightarrow\operatorname{Pic}(Y)
\]
is injective on geometric line-bundle classes if either
\begin{enumerate}
\item Frobenius is injective on $H^1(Y,\mathcal O_Y)$; or
\item the smooth variety $Y$ has a smooth lifting to $W_2(k)$.
\end{enumerate}
No lifting of $X$ is assumed in the second assertion, and no
lifting of the very ample line bundle on $Y$ is required.
The first assertion includes $H^1(Y,\mathcal O_Y)=0$.
Frobenius injectivity here is the precise cohomological
hypothesis; it is not replaced by a claim about the reduced
Picard variety alone.

Put $L=\mathcal O_Y(Y)$, which is very ample on $Y$.
If the point kernel were nonzero, the finite-$p$-group
reduction gives a nontrivial class of order $p$. By A2V,
there would then be an integer $a\geq1$ and a nonzero class
\[
 \alpha\in H^1(Y,L^{-a}),\qquad F\alpha=0
       \quad\hbox{in }H^1(Y,L^{-pa}).                  \tag{A2Y.1}
\]
We exclude this class under each hypothesis.

First suppose Frobenius is injective on $H^1(Y,\mathcal O_Y)$.
Since $\dim Y\geq2$ and $L^a$ is very ample, Bertini gives
a smooth connected divisor $E\in|L^a|$, cut out by a section
$s$. Projectivity and geometric connectedness give
$H^0(Y,\mathcal O_Y)=H^0(E,\mathcal O_E)=k$. The exact
sequence
\[
 0\longrightarrow L^{-a}\xrightarrow{\,s\,}\mathcal O_Y
  \longrightarrow\mathcal O_E\longrightarrow0
\]
therefore makes multiplication by $s$ an injection
\[
 j_s:H^1(Y,L^{-a})\hookrightarrow H^1(Y,\mathcal O_Y).
\]
The Frobenius square commutes:
\[
\begin{array}{ccc}
 H^1(Y,L^{-a})&\xrightarrow{\ F\ }&H^1(Y,L^{-pa})\\
 {\scriptstyle j_s}\downarrow&&\downarrow{\scriptstyle j_{s^p}}\\
 H^1(Y,\mathcal O_Y)&\xrightarrow{\ F\ }&H^1(Y,\mathcal O_Y).
\end{array}
\]
Its horizontal arrows are $p$-semilinear, or linear when
their targets are given the Frobenius scalar structure.
Equation (A2Y.1) implies $F(j_s\alpha)=0$. The assumed
injectivity of the bottom arrow and then the injectivity of
$j_s$ give $\alpha=0$, a contradiction. No assertion that
$j_{s^p}$ is injective is needed.

Now suppose that $Y$ lifts smoothly to $W_2(k)$. Apply
\cite[Section 11 hypothesis and Corollary 11.2, p.~2116]{DGJ_Langer2011}
with variety $Y$, numerically flat bundle $\mathcal O_Y$,
ample divisor corresponding to $L^a$, and indices $i=0,j=1$.
Its required inequality is
\[
 1<\min(p,\dim Y),
\]
which holds for every prime $p$ and every $\dim Y\geq2$.
It yields $H^1(Y,L^{-a})=0$, again contradicting (A2Y.1).
The proof of that corollary uses the truncated
Deligne--Illusie inequality only in total degree one and
iterates Frobenius until Serre vanishing applies. It does
not require $\dim Y<p$.

\emph{Remaining gate.}
A counterexample to point injectivity must therefore have
noninjective Frobenius on $H^1(Y,\mathcal O_Y)$ and a target
$Y$ that admits no smooth lifting to $W_2(k)$, as well as
the independent ambient restrictions of A2V and A2W. These
are necessary conditions, not an existence construction.
